% Trois médicaments et leurs groupes caractéristiques (exercice 8, chapitre 11)
\begin{tikzpicture}[x=1cm,y=1cm]
\node[anchor=west] at (0.1,0) {\chemfig[atom sep=1.7em]{HO-*6(=-=(-NH-[::+60](=[::+60]O)-[::-60]CH_3)-=-)}};
\node[pcLegendeG,anchor=north] at (2.1,-1.15) {paracétamol};
\node[anchor=west] at (6.1,0) {\chemfig[atom sep=1.7em]{*6(=-=(-O-[::+60](=[::+60]O)-[::-60]CH_3)-(-[::+60](=[::+60]O)-[::-60]OH)=-)}};
\node[pcLegendeG,anchor=north] at (7.9,-1.15) {aspirine};
\end{tikzpicture}
\par\vspace{2mm}
\begin{tikzpicture}[x=1cm,y=1cm]
\node[anchor=west] at (0.1,0) {\chemfig[atom sep=1.7em]{CH_3-CH(-[::-60]CH_3)-CH_2-*6(=-=(-CH(-[::+60]CH_3)-[::-60](=[::-120]O)-[::0]OH)-=-)}};
\node[pcLegendeG,anchor=north] at (3.6,-1.2) {ibuprofène};
\end{tikzpicture}
\par\smallskip\legende{Les formules de trois principes actifs courants.}
